\needspace{10\baselineskip}
\SourceLeaf{316}
\begin{center}
\pdfbookmark[1]{14. Verbesserung der Elemente anhand von vier Oppositionen}{zusatz-014}
[14.]
\end{center}
\begin{center}
\{Über die Verbesserung der Elemente eines Planeten\\anhand von vier beobachteten Oppositionen.\}
\end{center}
Sind die Elemente eines Planeten bereits so genau bekannt, dass die angenommenen Werte nur sehr wenig von den wahren abweichen, und liegen Beobachtungen von vier Oppositionen vor, lassen sich die Elemente leichter und einfacher verbessern als bei vier beliebig gewählten Planetenörtern. Die Theorie dieser Untersuchung wurde bereits in der Abhandlung »Disquisitio de elementis ellipticis Palladis« dargelegt. Hier folgen einige praktische Bemerkungen, die allgemein bei der Berechnung von Oppositionen und ähnlichen Fällen brauchbar sind.

Zuerst muss die Epoche verbessert werden. Man berechnet aus einer beliebigen Beobachtung, am besten aus einer, die zu einer der mittleren Oppositionen gehört, die geozentrische Länge und Breite des Planeten. Mithilfe des Radiusvektors, den man aus der unveränderten Epoche mit den bekannten Formeln erhält, berechnet man daraus die heliozentrische Länge und Breite und aus diesen die Länge in der Bahn. Nachdem man die Störungen bei dieser Länge berücksichtigt hat (sofern dafür bereits Formeln vorliegen), leitet man daraus die mittlere Anomalie ab. Vergleicht man sie mit der aus der unveränderten Epoche erhaltenen mittleren Anomalie, so ergibt sich die Änderung der Epoche. Liegen keine Beobachtungen vor und sind nur die berechneten Längen und Breiten zur Zeit der vier Oppositionen gegeben, so ist eine dieser Längen zu wählen. Es ist jedoch besser, die Beobachtungen selbst vorzunehmen, als sich auf die Berechnungen anderer zu verlassen.

Mit der so verbesserten Epoche und den übrigen, unveränderten Elementen berechnet man nun vorläufig die vier Oppositionen. Frühere Berechnungen lassen sich dabei mit hinreichender Genauigkeit verwenden. Zur möglichst einfachen Berechnung einer Opposition verfährt man folgendermaßen.

Liegen Beobachtungen von mehreren Orten vor, wählt man die Beobachtungen des Ortes, der die meisten geliefert hat. Aus den angegebenen Zeiten des Meridiandurchgangs bestimmt man die Kulminationszeiten des Planeten im Abstand von jeweils zwei Tagen für den gesamten Zeitraum, den die Beobachtungen abdecken. Durch eine vorläufige Rechnung für einen Tag in der Mitte der Beobachtungszeit bestimmt man die Entfernung des Planeten von der Erde und daraus die Zeit, die das Licht für den Weg vom Planeten zur Erde benötigt. Diese Zeit ist von den Kulminationszeiten abzuziehen. Damit ist die Aberration berücksichtigt. Aus der verbesserten Epoche leitet man nun durch Abziehen der Perihellänge zur Zeit der Epoche die mittlere Anomalie ab. Mithilfe der täglichen tropischen Bewegung bestimmt man daraus die mittleren Anomalien für die gewählten Tage und aus diesen die Rektaszensionen und Deklinationen. Um die Nutation zu berücksichtigen, verwendet man die wahre Länge des Knotens. Man erhält sie, indem man die Nutation in der Länge (\textsc{Bohnenbergers} Astron. p. 672) zur mittleren, durch die Elemente gegebenen Länge addiert. Die Parallaxe berücksichtigt man nach dem in der Theoria motus Corpor. Coel. art. 59 oder 70 erläuterten Verfahren. Es ist leicht einzusehen, warum die Rechnung auch dann im Abstand von zwei Tagen durchgeführt wird, wenn an einigen dieser Tage keine Beobachtungen vorliegen. Durch Differenzen kann man die Rechnung sehr bequem prüfen, sodass kein Fehler übersehen werden kann. Durch genaue Interpolation erhält man anschließend die Berechnungen für \textls[150]{alle} Tage.

Die so erhaltenen Rektaszensionen und Deklinationen vergleicht man mit den beobachteten Werten und bestimmt daraus die mittlere Abweichung sowohl der Rektaszension als auch der Deklination. Indem man diese Abweichungen zu den berechneten Werten addiert, leitet man die geozentrische Länge für den Tag der Opposition ab. Dieser Tag lässt sich sehr schnell \SourceLeaf{317}durch Vergleich der Beobachtungen mit den berechneten Sonnenlängen finden. Dieselbe Länge bestimmt man für einen Zeitpunkt, der um einen beliebigen Betrag (etwa einige Stunden) von der zuvor angenommenen Zeit abweicht. Nun sucht man durch Interpolation den Augenblick zu bestimmen, in dem sich die Längen der Sonne und des Planeten um genau \(180^{0}\) unterscheiden. Bei der Berechnung für die Sonne ist für die um die Aberration des Planeten korrigierte Kulminationszeit die \textls[150]{wahre} Länge zu suchen, also die anhand der Tafeln gefundene Länge, die die Aberration der Sonne bereits einschließt, um diese Aberration zu erhöhen.

Liegen nur wenige Beobachtungen vor (etwa bis zu sechs), so kann man die Rektaszensionen und Deklinationen auch nur für die Beobachtungstage berechnen, wenn das günstiger erscheint.

Bei der vorläufigen Berechnung der Oppositionen kann man die Aberration, Nutation und Parallaxe gegebenenfalls vernachlässigen. Will man die Oppositionen jedoch nur einmal berechnen, so muss man sie berücksichtigen.

Aus den für die Zeit der vier Oppositionen gefundenen Längen und Breiten leitet man nun nach dem in der Abhandlung de Elementis Palladis angegebenen Verfahren die Verbesserungen der einzelnen Elemente ab.

\par\medskip\begin{center}\rule{.22\textwidth}{.3pt}\end{center}
